\subsection{Fredholm alternative}

PROPOSITION 4 (Fredholm alternative). --- Let E be a Banach space and let h be a compact endomorphism of E. Let \(\lambda\) be an element of K-\{0\} . Let F be the kernel of \(1_{\mathrm{E}} - \lambda h\) and let G be the kernel of \(1_{\mathrm{E}'} - \lambda^{t}h\) .

\begin{enumerate}
\def\labelenumi{\alph{enumi})}
\item
  The spaces F and G have the same finite dimension n ≥ 0;
\item
  For \(y \in \mathrm{E}\) , there exists \(x \in \mathrm{E}\) such that \(x - \lambda h(x) = y\) if and only if \(\langle y, \ell \rangle = 0\) for all \(\ell \in \mathrm{G}\) ;
\item
  For \(\ell \in \mathrm{E}'\) , there exists \(f \in \mathrm{E}'\) such that \(f - \lambda^t h(f) = \ell\) if and only if \(\langle x, f \rangle = 0\) for all \(x \in \mathrm{F}\) .
\end{enumerate}

In particular, one and only one of the following conditions holds:

\begin{enumerate}
\def\labelenumi{(\roman{enumi})}
\item
  The space F is nonzero;
\item
  For every \(y \in \mathrm{E}\) , there exists \(x \in \mathrm{E}\) such that \(x - \lambda h(x) = y\) .
\end{enumerate}

By replacing h with \(\lambda h\) we reduce to the case where \(\lambda = 1\) .

The endomorphism \(^{t}h\) is compact by Corollary 1 of III, p. 9. The endomorphisms \(w = 1_{\mathrm{E}} - h\) of E and \(w' = 1_{\mathrm{E}'} - ^{t}h\) of \(\mathrm{E}'\) are therefore Riesz endomorphisms (cor. 2 of prop. 2 of III, p. 75). In particular, their kernels F and G are finite-dimensional. Since \(w' = ^{t}w\) and since the index of w is zero, the dimension of F is equal to that of G by formula (3) of III, p. 43.

For an element \(y\) of E, the equation \(x - h(x) = y\) has a solution \(x \in \mathrm{E}\) if and only if \(x\) belongs to the image of \(w\) . Since the dual of the cokernel of \(w\) is identified with the kernel G of \(\operatorname{Ker}(^t w)\) (EVT, IV, p. 27, prop. 2), this is the case if and only if \(\langle y, \ell \rangle = 0\) for all \(\ell \in \mathrm{G}\) .

For an element \(\ell\) of E', the equation \(f - {}^t h(f) = \ell\) has a solution \(f \in \mathrm{E}'\) if and only if \(f\) belongs to the image of \(w'\) . Since the image of \({}^t w\) is the orthogonal of the kernel F of \(w\) (EVT, IV, p. 27, prop. 2) , this is the case if and only if \(\langle x, f \rangle = 0\) for all \(x \in \mathrm{F}\) .

Example. --- Let X be a compact space, \(\mu\) a measure on X, real or complex depending on whether K is equal to R or C, and N: \(X \times X \to K\) a continuous function.

Let E denote the Banach space \(\mathcal{C}(\mathrm{X};\mathrm{K})\) . Suppose that a function \(a\in\mathrm{E}\) and \(\lambda\in\mathrm{K}-\{0\}\) are given. We study the problem of the existence and uniqueness of solutions \(f\in\mathrm{E}\) of the integral equation

\[
f (x) - \lambda \int_ {\mathrm{X}} \mathrm{N} (x, y) f (y) d \mu (y) = a (x) \quad (x \in \mathrm{X}).\tag{3}
\]

For this, let us introduce the set \(F_{\lambda}\) of solutions \(f \in E\) of the associated homogeneous equation

\[
f (x) - \lambda \int_ {\mathrm{X}} \mathrm{N} (x, y) f (y) d \mu (y) = 0 \qquad (x \in \mathrm{X}),\tag{4}
\]

and the set \(G_{\lambda}\) of solutions \(g \in E\) of the homogeneous transposed equation of (4), that is,

\[
g (y) - \lambda \int_ {\mathrm{X}} \mathrm{N} (x, y) g (x) d \mu (x) = 0 \qquad (y \in \mathrm{X}).\tag{5}
\]

If it is not empty, the set of solutions of (3) is an affine subspace of E with direction \(F_{\lambda}\) .

For \(f\in \mathrm{E}\) and \(x\in \mathrm{X}\) , set

\[
h (f) (x) = \int_ {\mathrm{X}} \mathrm{N} (x, y) f (y) d \mu (y);
\]

the map \(f \mapsto h(f)\) thus defined is a compact endomorphism of the Banach space E (cor. of prop. 2 of III, p. 26). The dual E' of E consists of measures on X, real or complex depending on whether K equals R or C (INT, III, p. 47, § 1, n° 3, def. 2). The transpose \(^{t}h\) of h is the endomorphism of E' characterized by the relation \(\langle f, ^{t}h(m)\rangle = \langle h(f), m\rangle\) for \(f \in E\) and \(m \in E'\) , that is,

\[
\begin{array}{r l} \int_ {\mathrm{X}} f (x) d (^ {t} h (m)) (x) & = \int_ {\mathrm{X}} h (f) (x) d m (x) \\ & = \int_ {\mathrm{X}} \Big (\int_ {\mathrm{X}} \mathrm{N} (x, y) f (y) d \mu (y) \Big) d m (x) \\ & = \int_ {\mathrm{X}} \Big (\int_ {\mathrm{X}} \mathrm{N} (x, y) d m (x) \Big) f (y) d \mu (y) \end{array}\tag{6}
\]

for \(m \in E'\) and \(f \in E\) (INT, III, p. 84, § 4, no. 1, th. 2).

Lemma 4. --- The linear map from \(G_{\lambda}\) into \(E'\) defined by \(g \mapsto g \cdot \mu\) is injective, and its image is the kernel of \(1_{E'} - \lambda^{t}h\) .

If \(g \in G_{\lambda}\) satisfies \(g \cdot \mu = 0\) then formula (5) implies \(g = 0\) , so the map \(g \mapsto g \cdot \mu\) is injective.

Let \(m \in \mathrm{E}'\) be a measure belonging to the kernel of \(1_{\mathrm{E}'} - \lambda^t h\) . By relation (6), we have

\[
\int_ {\mathrm{X}} f (x) d m (x) = \lambda \int_ {\mathrm{X}} \left(\int_ {\mathrm{X}} \mathrm{N} (x, y) d m (x)\right) f (y) d \mu (y)\tag{7}
\]

for every continuous function \(f: \mathrm{X} \to \mathrm{K}\) . The measure \(m\) is therefore equal to the measure \(g \cdot \mu\) , where \(g\) is the continuous function from \(\mathrm{X}\) into \(\mathrm{K}\) defined by

\[
g (y) = \lambda \int_ {\mathrm{X}} \mathrm{N} (x, y) d m (x)\tag{8}
\]

for all \(y \in X\) . Formula (8) then implies that the function g belongs to \(G_{\lambda}\) .

Conversely, let \(g: X \to K\) be a continuous function belonging to \(G_{\lambda}\) and set \(m = g \cdot \mu\) . For every continuous function \(f: X \to K\) , we therefore have

\[
\begin{array}{l} \int_ {\mathrm{X}} f (y) d m (y) = \int_ {\mathrm{X}} g (y) f (y) d \mu (y) \\ \qquad = \lambda \int_ {\mathrm{X}} \Big (\int_ {\mathrm{X}} \mathrm{N} (x, y) g (x) d \mu (x) \Big) f (y) d \mu (y) \\ \qquad = \lambda \int_ {\mathrm{X}} f (x) d m (x), \end{array}
\]

so that \({}^{t}h(m)=\lambda m\) . This proves Lemma 4.

Applying Proposition 4 and Theorem 4 of III, p. 78 to the endomorphism h, we then obtain the following statements.

THEOREM 5. --- Let \(\lambda\in K\) .

\begin{enumerate}
\def\labelenumi{\alph{enumi})}
\item
  The sets \(F_{\lambda}\) and \(G_{\lambda}\) are finite-dimensional vector subspaces of \(\mathcal{C}(X;K)\) and their dimensions are equal;
\item
  For equation (3) to have at least one solution \(f \in \mathcal{C}(\mathrm{X}, \mathrm{K})\) it is necessary and sufficient that one has \(\int_{\mathrm{X}} a(x) g(x) d\mu(x) = 0\) for every function \(g \in G_{\lambda}\) . The set of solutions of (3) is then an affine subspace of \(\mathcal{C}(\mathrm{X}; \mathrm{K})\) with direction \(F_{\lambda}\) ;
\item
  One of the following mutually exclusive conditions is satisfied:
\end{enumerate}

\begin{enumerate}
\def\labelenumi{(\roman{enumi})}
\item
  For any function \(a \in \mathcal{C}(\mathrm{X};\mathrm{K})\) there exists a unique solution \(f \in \mathcal{C}(\mathrm{X};\mathrm{K})\) of the integral equation (3);
\item
  The homogeneous equation (4) has a nonzero solution \(f \in \mathcal{C}(\mathrm{X};\mathrm{K})\) .
\end{enumerate}

Corollary. --- The set of \(\lambda\in\mathrm{K}\) for which the space \(\mathrm{F}_{\lambda}\) is nonzero is a countable, closed, and discrete subset of K.

